\subsection{Test di accettazione}
\label{subsec:test_accettazione}
I test di accettazione verificano che il sistema soddisfi i requisiti delineati nel documento di Analisi dei Requisiti e le aspettative degli utenti finali.
Vengono eseguiti dagli utenti finali in scenari realistici per assicurarsi che tutte le funzionalità operino come previsto prima del rilascio ufficiale.

\begin{tabularx}{\textwidth}{|c|X|c|}
    \hline
    \textbf{ID} & \textbf{Descrizione} & \textbf{Stato}\\
    \hline
    \label{ta:ta1}T.A.1 & Verificare che l'utente possa visualizzare il dataset caricato:
                        \begin{itemize}
                            \item visualizzare lo stato del dataset (se temporaneo o se salvato);
                            \item visualizzare il nome del dataset;
                            \item visualizzare la lista di coppie domanda e risposta del dataset.
                        \end{itemize} & Non superato\\ %uc1
    \hline
    \label{ta:ta2}T.A.2 & Verificare che l'utente possa effettuare una ricerca tra gli elementi di un dataset caricato:
                        \begin{itemize}
                            \item visualizzare la lista di coppie domanda e risposta del dataset;
                            \item richiedere al sistema di effettuare una ricerca data una stringa;
                            \item visualizzare le coppie domanda e risposta che contengono la stringa.
                        \end{itemize} & Non superato\\%uc4
    \hline
    \label{ta:ta3}T.A.3 & Verificare che l'utente possa modificare gli elementi del dataset visualizzato:
                        \begin{itemize}
                            \item visualizzare la lista di coppie domanda e risposta del dataset;
                            \item eliminare, modificare o aggiungere una coppia domanda e risposta;
                            \item confermare la modifica attuata.
                        \end{itemize} & Non superato\\%uc5 uc6 uc7 
    \hline
    \label{ta:ta4}T.A.4 & Verificare che l'utente possa salvare il dataset:
                        \begin{itemize}
                            \item essere nella pagina di visualizzazione del dataset;
                            \item assegnare un nome al dataset se è nuovo;
                            \item scegliere se creare una nuova versione aggiornata del dataset o se eliminare i test associati al dataset, se il dataset era già stato salvato in precedenza;
                            \item confermare il salvataggio del dataset.
                        \end{itemize} & Superato\\ %uc11 uc12 uc14
    \hline
    \label{ta:ta5}T.A.5 & Verificare che l'utente possa visualizzare una lista dei dataset salvati:
                        \begin{itemize}
                            \item richiedere la visualizzazione dei dataset salvati;
                            \item visualizzare i nomi dei dataset salvati;
                            \item visualizzare la data dell'ultima modifica per ogni dataset.
                        \end{itemize} & Superato\\ %uc19
    \hline
    \label{ta:ta6}T.A.6 & Verificare che l'utente possa caricare un dataset da un file JSON:
                        \begin{itemize}
                            \item visualizzare la lista dei dataset salvati;
                            \item richiedere il salvataggio di un nuovo dataset a partire da un file JSON;
                            \item selezionare il file JSON.
                        \end{itemize} & Non superato\\ %uc17
    \hline
    \label{ta:ta7}T.A.7 & Verificare che l'utente possa effettuare una ricerca per nome tra i dataset salvati:
                        \begin{itemize}
                            \item visualizzare la lista dei dataset salvati;
                            \item richiedere al sistema di effettuare una ricerca data una stringa;
                            \item ottenere i risultati della ricerca.
                        \end{itemize} & Superato\\ %uc20
    \hline
    \label{ta:ta8}T.A.8 & Verificare che l'utente possa gestire un dataset salvato:
                        \begin{itemize}
                            \item visualizzare la lista dei dataset salvati;
                            \item rinomimare un dataset salvato;
                            \item copiare un dataset salvato;
                            \item confermare l'operazione di gestione richiesta.
                        \end{itemize} & Superato\\ %uc24 uc21
    \hline
    \label{ta:ta9}T.A.9 & Verificare che l'utente possa eliminare un dataset salvato:
                        \begin{itemize}
                            \item visualizzare la lista dei dataset salvati;
                            \item selezionare il dataset da eliminare;
                            \item visualizzare la lista dei test associati al dataset che verranno eliminati;
                            \item confermare l'operazione di eliminazione.
                        \end{itemize} & Superato\\ %uc22 uc23
    \hline
    \label{ta:ta10}T.A.10 & Verificare che l'utente possa cambiare il dataset caricato:
                        \begin{itemize}
                            \item visualizzare la lista dei dataset salvati;
                            \item caricare un nuovo dataset vuoto temporaneo;
                            \item caricare un dataset salvato;
                            \item decidere se sovrascrivere il dataset caricato attuale.
                        \end{itemize} & Non superato\\ %uc25 uc26 uc27
    \hline
    \label{ta:ta11}T.A.11 & Verificare che l'utente possa eseguire il test sul dataset caricato:
                        \begin{itemize}
                            \item richiedere di eseguire il test;
                            \item selezionare l'LLM da testare tra quelli salvati;
                            \item visualizzare i risultati del test eseguito.
                        \end{itemize} & Non superato\\ %uc28 uc31 uc32 uc33 uc34
    \hline
    \label{ta:ta12}T.A.12 & Verificare che l'utente possa salvare i risultati di un test eseguito:
                        \begin{itemize}
                            \item essere nella pagina di visualizzazione dei risultati del test;
                            \item assegnare un nome al test;
                            \item confermare il salvataggio.
                        \end{itemize} & Non superato\\ %uc35
    \hline
    \label{ta:ta13}T.A.13 & Verificare che l'utente possa visualizzare una lista dei test salvati:
                        \begin{itemize}
                            \item richiedere la visualizzazione dei test salvati;
                            \item visualizzare i nomi dei test salvati;
                            \item visualizzare, per ogni test, il nome del dataset su cui è stato eseguito e la data di esecuzione.
                        \end{itemize} & Non superato\\ %uc36
    \hline
    \label{ta:ta14}T.A.14 & Verificare che l'utente possa effettuare una ricerca per nome tra i test salvati:
                        \begin{itemize}
                            \item visualizzare la lista dei test salvati;
                            \item richiedere al sistema di effettuare una ricerca data una stringa;
                            \item ottenere i risultati della ricerca.
                        \end{itemize} & Non superato\\ %uc37
    \hline
    \label{ta:ta15}T.A.15 & Verificare che l'utente possa cambiare il test caricato:
                        \begin{itemize}
                            \item visualizzare la lista dei test salvati;
                            \item selezionare il test salvato da caricare;
                            \item decidere se sovrascrivere il test caricato attuale;
                            \item visualizzare i risultati del nuovo test caricato.
                        \end{itemize} & Non superato\\ %uc38 uc39
    \hline
    \label{ta:ta16}T.A.16 & Verificare che l'utente possa gestire un test salvato:
                        \begin{itemize}
                            \item visualizzare la lista dei test salvati;
                            \item rinomimare un test salvato;
                            \item eliminare un test salvato;
                            \item confermare l'operazione richiesta.
                        \end{itemize} & Non superato\\%uc40 uc41
    \hline
    \label{ta:ta17}T.A.17 & Verificare che l'utente possa confrontare il test caricato con uno dei test salvati effettuati sullo stesso dataset o LLM:
                        \begin{itemize}
                            \item richiedere l'esecuzione del confronto;
                            \item selezionare il test salvato da confrontare;
                            \item visualizzare il confronto tra i due test.
                        \end{itemize} & Non superato\\ %uc42 uc43 uc44 uc45
    \hline
    \label{ta:ta18}T.A.18 & Verificare che l'utente possa visualizzare una lista degli LLM da testare salvati:
                        \begin{itemize}
                            \item richiedere la visualizzazione degli LLM salvati;
                            \item visualizzare i nomi degli LLM salvati.
                        \end{itemize} & Non superato\\ %uc52
    \hline
    \label{ta:ta19}T.A.19 & Verificare che l'utente possa visualizzare un singolo LLM salvato:
                        \begin{itemize}
                            \item visualizzare la lista degli LLM salvati;
                            \item selezionare un LLM da visualizzare;
                            \item visualizzare il nome e la configurazione dell'LLM selezionato.
                        \end{itemize} & Non superato\\ %uc53
    \hline
    \label{ta:ta20}T.A.20 & Verificare che l'utente possa eliminare un LLM salvato:
                        \begin{itemize}
                            \item visualizzare la lista degli LLM salvati;
                            \item selezionare l'LLM salvato da eliminare;
                            \item visualizzare la lista dei test associati all'LLM che verranno eliminati;
                            \item confermare l'eliminazione dell'LLM.
                        \end{itemize} & Non superato\\ %uc54 uc55
    \hline
    \label{ta:ta21}T.A.21 & Verificare che l'utente possa aggiungere un nuovo LLM:
                        \begin{itemize}
                            \item visualizzare la lista degli LLM salvati;
                            \item richiedere di aggiungere un nuovo LLM;
                            \item assegnare un nome al nuovo LLM;
                            \item specificare la configurazione delle richieste HTTP al nuovo LLM;
                            \item salvare il nuovo LLM.
                        \end{itemize} & Non superato\\ %uc46 uc47
    \hline
    \label{ta:ta22}T.A.22 & Verificare che l'utente possa modificare un LLM salvato:
                        \begin{itemize}
                            \item visualizzare la lista degli LLM salvati;
                            \item richiedere la modifica di un LLM salvato;
                            \item rinominare l'LLM di cui si è richiesta la modifica;
                            \item modificare la configurazione delle richieste HTTP dell'LLM di cui si è richiesta la modifica (modificando, aggiungendo o eliminando elementi);
                            \item confermare la modifica attuata.
                        \end{itemize} & Non superato\\ %uc50 uc51
    \hline
\end{tabularx}

\begin{tabularx}{\textwidth}{|c|X|}
    \hline
    \textbf{ID} & \textbf{Caso d'uso}\\
    \hline
    \label{ta:trace_ta1}T.A.1 & UC-1\\
    \hline
    \label{ta:trace_ta2}T.A.2 & UC-4\\
    \hline
    \label{ta:trace_ta3}T.A.3 & UC-5, UC-6, UC-7\\
    \hline
    \label{ta:trace_ta4}T.A.4 & UC-11, UC-12, UC-14\\
    \hline
    \label{ta:trace_ta5}T.A.5 & UC-19\\
    \hline
    \label{ta:trace_ta6}T.A.6 & UC-17\\
    \hline
    \label{ta:tarce_ta7}T.A.7 & UC-20\\
    \hline
    \label{ta:trace_ta8}T.A.8 & UC-21, UC-24\\
    \hline
    \label{ta:trace_ta9}T.A.9 & UC-22, UC-23\\
    \hline
    \label{ta:trace_ta10}T.A.10 & UC-25, UC-26, UC-27\\
    \hline
    \label{ta:trace_ta11}T.A.11 & UC-28, UC-31, UC-32, UC-33, UC-34\\
    \hline
    \label{ta:trace_ta12}T.A.12 & UC-35\\
    \hline
    \label{ta:trace_ta13}T.A.13 & UC-36\\
    \hline
    \label{ta:trace_ta14}T.A.14 & UC-37\\
    \hline
    \label{ta:trace_ta15}T.A.15 & UC-38, UC-39\\
    \hline
    \label{ta:trace_ta16}T.A.16 & UC-40, UC-41\\
    \hline
    \label{ta:trace_ta17}T.A.17 & UC-42, UC-43, UC-44, UC-45\\
    \hline 
    \label{ta:trace_ta18}T.A.18 & UC-52\\
    \hline
    \label{ta:trace_ta19}T.A.19 & UC-53\\
    \hline
    \label{ta:trace_ta20}T.A.20 & UC-54, UC-55\\
    \hline
    \label{ta:trace_ta21}T.A.21 & UC-46, UC-47\\
    \hline
    \label{ta:tarce_ta22}T.A.22 & UC-50, UC-51\\
    \hline
\end{tabularx}




